% 原创TikZ重绘；层次布局参考Bishop与Bishop（2024）第12章。
\begin{tikzpicture}[x=1mm,y=1mm,font=\figurefont\footnotesize,
 every node/.style={text=NNInk},
 stage/.style={nn operator,minimum height=9mm,inner sep=2pt},
 norm/.style={stage,minimum height=7mm},
 wire/.style={nn flow,draw=NNWire,rounded corners=1.5mm},
 aux/.style={draw=NNLine,line width=.85pt,densely dashed},
 sum/.style={nn pointwise,minimum size=5mm}]

\node[] (stoken0) at (17,0) {猫};
\node[inner sep=0pt,minimum width=12mm,minimum height=3mm] (sembed0) at (17,13) {};
\filldraw[draw=NNInput,fill=NNInput!18,line width=.5pt] (11.0,11.5) rectangle ++(3.0,3);
\filldraw[draw=NNInput,fill=NNInput!18,line width=.5pt] (14.0,11.5) rectangle ++(3.0,3);
\filldraw[draw=NNInput,fill=NNInput!18,line width=.5pt] (17.0,11.5) rectangle ++(3.0,3);
\filldraw[draw=NNInput,fill=NNInput!18,line width=.5pt] (20.0,11.5) rectangle ++(3.0,3);
\draw[wire] (stoken0.north)--(sembed0.south);
\node[sum,draw=NNInput] (sadd0) at (17,29) {$+$};
\draw[wire] (sembed0.north)--(sadd0.south);
\node[nn annotation] (spe0) at (10,39) {$\bm p_{0}$};
\draw[wire] (spe0.south)--(sadd0.north west);
\draw[wire] (sadd0.north)--(17,51);
\node[] (stoken1) at (35,0) {追};
\node[inner sep=0pt,minimum width=12mm,minimum height=3mm] (sembed1) at (35,13) {};
\filldraw[draw=NNInput,fill=NNInput!18,line width=.5pt] (29.0,11.5) rectangle ++(3.0,3);
\filldraw[draw=NNInput,fill=NNInput!18,line width=.5pt] (32.0,11.5) rectangle ++(3.0,3);
\filldraw[draw=NNInput,fill=NNInput!18,line width=.5pt] (35.0,11.5) rectangle ++(3.0,3);
\filldraw[draw=NNInput,fill=NNInput!18,line width=.5pt] (38.0,11.5) rectangle ++(3.0,3);
\draw[wire] (stoken1.north)--(sembed1.south);
\node[sum,draw=NNInput] (sadd1) at (35,29) {$+$};
\draw[wire] (sembed1.north)--(sadd1.south);
\node[nn annotation] (spe1) at (28,39) {$\bm p_{1}$};
\draw[wire] (spe1.south)--(sadd1.north west);
\draw[wire] (sadd1.north)--(35,51);
\node[] (stoken2) at (53,0) {狗};
\node[inner sep=0pt,minimum width=12mm,minimum height=3mm] (sembed2) at (53,13) {};
\filldraw[draw=NNInput,fill=NNInput!18,line width=.5pt] (47.0,11.5) rectangle ++(3.0,3);
\filldraw[draw=NNInput,fill=NNInput!18,line width=.5pt] (50.0,11.5) rectangle ++(3.0,3);
\filldraw[draw=NNInput,fill=NNInput!18,line width=.5pt] (53.0,11.5) rectangle ++(3.0,3);
\filldraw[draw=NNInput,fill=NNInput!18,line width=.5pt] (56.0,11.5) rectangle ++(3.0,3);
\draw[wire] (stoken2.north)--(sembed2.south);
\node[sum,draw=NNInput] (sadd2) at (53,29) {$+$};
\draw[wire] (sembed2.north)--(sadd2.south);
\node[nn annotation] (spe2) at (46,39) {$\bm p_{2}$};
\draw[wire] (spe2.south)--(sadd2.north west);
\draw[wire] (sadd2.north)--(53,51);
\node[stage,minimum width=58mm,minimum height=10mm] (slayer1) at (35,57) {第$1$层：双向注意力＋FFN};
\node[stage,minimum width=58mm,minimum height=10mm] (slayerL) at (35,89) {第$L$层：双向注意力＋FFN};
\draw[wire] (17,62)--(17,67);
\node[] (sdots17) at (17,73) {$\vdots$};
\draw[wire] (17,79)--(17,84);
\draw[wire] (35,62)--(35,67);
\node[] (sdots35) at (35,73) {$\vdots$};
\draw[wire] (35,79)--(35,84);
\draw[wire] (53,62)--(53,67);
\node[] (sdots53) at (53,73) {$\vdots$};
\draw[wire] (53,79)--(53,84);
\node[nn annotation,anchor=west] at (61,73) {$\times L$};
\node[] (ttoken0) at (104,0) {BOS};
\node[inner sep=0pt,minimum width=12mm,minimum height=3mm] (tembed0) at (104,13) {};
\filldraw[draw=NNInput,fill=NNInput!18,line width=.5pt] (98.0,11.5) rectangle ++(3.0,3);
\filldraw[draw=NNInput,fill=NNInput!18,line width=.5pt] (101.0,11.5) rectangle ++(3.0,3);
\filldraw[draw=NNInput,fill=NNInput!18,line width=.5pt] (104.0,11.5) rectangle ++(3.0,3);
\filldraw[draw=NNInput,fill=NNInput!18,line width=.5pt] (107.0,11.5) rectangle ++(3.0,3);
\draw[wire] (ttoken0.north)--(tembed0.south);
\node[sum,draw=NNInput] (tadd0) at (104,29) {$+$};
\draw[wire] (tembed0.north)--(tadd0.south);
\node[nn annotation] (tpe0) at (97,39) {$\bm p_{0}$};
\draw[wire] (tpe0.south)--(tadd0.north west);
\draw[wire] (tadd0.north)--(104,51);
\node[] (ttoken1) at (121,0) {cat};
\node[inner sep=0pt,minimum width=12mm,minimum height=3mm] (tembed1) at (121,13) {};
\filldraw[draw=NNInput,fill=NNInput!18,line width=.5pt] (115.0,11.5) rectangle ++(3.0,3);
\filldraw[draw=NNInput,fill=NNInput!18,line width=.5pt] (118.0,11.5) rectangle ++(3.0,3);
\filldraw[draw=NNInput,fill=NNInput!18,line width=.5pt] (121.0,11.5) rectangle ++(3.0,3);
\filldraw[draw=NNInput,fill=NNInput!18,line width=.5pt] (124.0,11.5) rectangle ++(3.0,3);
\draw[wire] (ttoken1.north)--(tembed1.south);
\node[sum,draw=NNInput] (tadd1) at (121,29) {$+$};
\draw[wire] (tembed1.north)--(tadd1.south);
\node[nn annotation] (tpe1) at (114,39) {$\bm p_{1}$};
\draw[wire] (tpe1.south)--(tadd1.north west);
\draw[wire] (tadd1.north)--(121,51);
\node[] (ttoken2) at (138,0) {chases};
\node[inner sep=0pt,minimum width=12mm,minimum height=3mm] (tembed2) at (138,13) {};
\filldraw[draw=NNInput,fill=NNInput!18,line width=.5pt] (132.0,11.5) rectangle ++(3.0,3);
\filldraw[draw=NNInput,fill=NNInput!18,line width=.5pt] (135.0,11.5) rectangle ++(3.0,3);
\filldraw[draw=NNInput,fill=NNInput!18,line width=.5pt] (138.0,11.5) rectangle ++(3.0,3);
\filldraw[draw=NNInput,fill=NNInput!18,line width=.5pt] (141.0,11.5) rectangle ++(3.0,3);
\draw[wire] (ttoken2.north)--(tembed2.south);
\node[sum,draw=NNInput] (tadd2) at (138,29) {$+$};
\draw[wire] (tembed2.north)--(tadd2.south);
\node[nn annotation] (tpe2) at (131,39) {$\bm p_{2}$};
\draw[wire] (tpe2.south)--(tadd2.north west);
\draw[wire] (tadd2.north)--(138,51);
\node[] (ttoken3) at (155,0) {dog};
\node[inner sep=0pt,minimum width=12mm,minimum height=3mm] (tembed3) at (155,13) {};
\filldraw[draw=NNInput,fill=NNInput!18,line width=.5pt] (149.0,11.5) rectangle ++(3.0,3);
\filldraw[draw=NNInput,fill=NNInput!18,line width=.5pt] (152.0,11.5) rectangle ++(3.0,3);
\filldraw[draw=NNInput,fill=NNInput!18,line width=.5pt] (155.0,11.5) rectangle ++(3.0,3);
\filldraw[draw=NNInput,fill=NNInput!18,line width=.5pt] (158.0,11.5) rectangle ++(3.0,3);
\draw[wire] (ttoken3.north)--(tembed3.south);
\node[sum,draw=NNInput] (tadd3) at (155,29) {$+$};
\draw[wire] (tembed3.north)--(tadd3.south);
\node[nn annotation] (tpe3) at (148,39) {$\bm p_{3}$};
\draw[wire] (tpe3.south)--(tadd3.north west);
\draw[wire] (tadd3.north)--(155,51);
\node[stage,minimum width=66mm,minimum height=10mm] (tlayer1) at (129,57) {第$1$层：因果＋交叉注意力＋FFN};
\node[stage,minimum width=66mm,minimum height=10mm] (tlayerL) at (129,89) {第$L$层：因果＋交叉注意力＋FFN};
\draw[wire] (104,62)--(104,67);
\node[] (tdots104) at (104,73) {$\vdots$};
\draw[wire] (104,79)--(104,84);
\draw[wire] (121,62)--(121,67);
\node[] (tdots121) at (121,73) {$\vdots$};
\draw[wire] (121,79)--(121,84);
\draw[wire] (138,62)--(138,67);
\node[] (tdots138) at (138,73) {$\vdots$};
\draw[wire] (138,79)--(138,84);
\draw[wire] (155,62)--(155,67);
\node[] (tdots155) at (155,73) {$\vdots$};
\draw[wire] (155,79)--(155,84);
\node[nn annotation,anchor=west] at (163,73) {$\times L$};
\draw[wire] (17,94)--(17,107);
\node[inner sep=0pt,minimum width=12mm,minimum height=3mm] (sh0) at (17,109) {};
\filldraw[draw=NNOutput,fill=NNOutput!18,line width=.5pt] (11.0,107.5) rectangle ++(3.0,3);
\filldraw[draw=NNOutput,fill=NNOutput!18,line width=.5pt] (14.0,107.5) rectangle ++(3.0,3);
\filldraw[draw=NNOutput,fill=NNOutput!18,line width=.5pt] (17.0,107.5) rectangle ++(3.0,3);
\filldraw[draw=NNOutput,fill=NNOutput!18,line width=.5pt] (20.0,107.5) rectangle ++(3.0,3);
\draw[wire] (35,94)--(35,107);
\node[inner sep=0pt,minimum width=12mm,minimum height=3mm] (sh1) at (35,109) {};
\filldraw[draw=NNOutput,fill=NNOutput!18,line width=.5pt] (29.0,107.5) rectangle ++(3.0,3);
\filldraw[draw=NNOutput,fill=NNOutput!18,line width=.5pt] (32.0,107.5) rectangle ++(3.0,3);
\filldraw[draw=NNOutput,fill=NNOutput!18,line width=.5pt] (35.0,107.5) rectangle ++(3.0,3);
\filldraw[draw=NNOutput,fill=NNOutput!18,line width=.5pt] (38.0,107.5) rectangle ++(3.0,3);
\draw[wire] (53,94)--(53,107);
\node[inner sep=0pt,minimum width=12mm,minimum height=3mm] (sh2) at (53,109) {};
\filldraw[draw=NNOutput,fill=NNOutput!18,line width=.5pt] (47.0,107.5) rectangle ++(3.0,3);
\filldraw[draw=NNOutput,fill=NNOutput!18,line width=.5pt] (50.0,107.5) rectangle ++(3.0,3);
\filldraw[draw=NNOutput,fill=NNOutput!18,line width=.5pt] (53.0,107.5) rectangle ++(3.0,3);
\filldraw[draw=NNOutput,fill=NNOutput!18,line width=.5pt] (56.0,107.5) rectangle ++(3.0,3);
\node[nn heading] (hs) at (35,121) {编码器最终输出 $\bm H_{\mathrm s}$};
\draw[draw=NNWire,line width=1.2pt] (17,115)--(79,115)--(79,57);
\draw[draw=NNWire,line width=1.2pt] (17,110.5)--(17,115);
\draw[draw=NNWire,line width=1.2pt] (35,110.5)--(35,115);
\draw[draw=NNWire,line width=1.2pt] (53,110.5)--(53,115);
\node[rotate=90,fill=none,inner sep=1pt,text=FigureInk] at (75,86)
  {KV bus};
\draw[wire] (79,57)--(tlayer1.west);
\fill[NNWire] (79,57) circle (.55);
\node[nn annotation] (kv1) at (87,63) {K、V};
\draw[wire] (79,89)--(tlayerL.west);
\fill[NNWire] (79,89) circle (.55);
\node[nn annotation] (kvL) at (87,95) {K、V};
\node[stage,draw=NNOutput,fill=NNOutput!18,minimum width=14mm,minimum height=8mm] (lm0) at (104,108) {LM};
\draw[wire] (104,94)--(lm0.south);
\node[] (prediction0) at (104,124) {cat};
\draw[wire] (lm0.north)--(prediction0.south);
\node[stage,draw=NNOutput,fill=NNOutput!18,minimum width=14mm,minimum height=8mm] (lm1) at (121,108) {LM};
\draw[wire] (121,94)--(lm1.south);
\node[] (prediction1) at (121,124) {chases};
\draw[wire] (lm1.north)--(prediction1.south);
\node[stage,draw=NNOutput,fill=NNOutput!18,minimum width=14mm,minimum height=8mm] (lm2) at (138,108) {LM};
\draw[wire] (138,94)--(lm2.south);
\node[] (prediction2) at (138,124) {dog};
\draw[wire] (lm2.north)--(prediction2.south);
\node[stage,draw=NNOutput,fill=NNOutput!18,minimum width=14mm,minimum height=8mm] (lm3) at (155,108) {LM};
\draw[wire] (155,94)--(lm3.south);
\node[] (prediction3) at (155,124) {EOS};
\draw[wire] (lm3.north)--(prediction3.south);
\node[nn heading] (enc) at (35,141) {编码器：读取全部源词元};
\node[nn heading] (dec) at (129,141) {解码器：读取已知目标前缀};
\node[nn annotation] (inputnote) at (35,-13) {源输入};
\node[nn annotation] (targetnote) at (129,-13) {右移后的目标输入};
\end{tikzpicture}
